\clearpage
\section{\label{sec:collection}Collecting Infrastructure Code as a Data Source}

\justifying
Some level of familiarity on the reader's part with \href{https://www.terraform.io/intro/index.html}{Terraform} and
by extension, \href{https://learn.hashicorp.com/tutorials/terraform/infrastructure-as-code}{Infrastructure as Code} is assumed.
\vspace{2mm}

\justifying
Terraform users create files that define network infrastructure. When this infrastructure is deployed via Terraform, a state file is created to keep track of the current deployment. This state file is essentially graph structured data. We can export this graph data from the actual (or planned) state and ingest this
with our deep learning pipeline.

\subsection{\label{sec:collect}The Data Collection Process}

\justifying
A tool for gathering and uploading from data sources is used to send relevant data from various sources to a storage bucket in
Google Cloud. Raw data is processed and tagged with a UUID as a unique key to associate data with it's source. The source code for this Python tool is 
\href{https://github.com/devsecfranklin/model-graph-neural-net/tree/main/src}{available in the Github repository for this project}.
\vspace{2mm}

\justifying
Identifying and versioning data as we would any other code or documentation artifact eases the
burden of handling our data and allows us to more easily work with larger data sets. While the collection process is relatively straight
forward to implement and maintain for this project on a small hobbyist/researcher scale, going forward it would make more sense to use
a \href{https://github.com/iterative/dvc}{community supported framework such as dvc} to leverage the efforts of folks doing similar work.

\section{\label{sec:iacdata}Preparing Infrastructure Code as Data}

\justifying
Once a Terraform code base has been initialized, we generate a directed graph of the infrastructure using Dot.
Dot is \href{https://graphviz.org/doc/info/lang.html}{a language used to describe graphs}. Consider the following
declaration of a Cloud Function for Google Cloud.

\begin{mybox}{\thetcbcounter: Terraform Declaration Example}
	\lstinputlisting{code/main.tf}
\end{mybox}

\justifying
\href{https://graphviz.org/}{Graphviz is open source graph visualization software}. Graph visualization is a way
of representing structural information as diagrams of abstract graphs and networks.
Conversion of Terraform HCL to a directed graph representation is trivial thanks to Graphviz and Dot. It is possible to reproduce the
digraph displayed in figure \ref{pygraph} \href{https://github.com/devsecfranklin/model-graph-neural-net/blob/main/bin/gen_graph.sh}{using the shell script provided}.

\justifying
\begin{figure}[H]
	\includegraphics[width=12cm]{graph.png}
	\caption{DOT output from Terraform Files as a digraph}
	\label{pygraph}
\end{figure}

\section{\label{sec:collection}Data Storage}

\justifying
Run a small tool to pass results to a Google Cloud storage bucket. Perhaps the DVC tool can assist with this part.
